% migrate-to: options-execution
\subsection{Options Market Microstructure}

\begin{principle} \hlt{Options Are Surface Markets}\\
An option quote is linked to the underlying price, volatility, rates, dividends, strike, and maturity.\\
Liquidity therefore lives on a surface, not in one instrument, and market makers hedge across the underlying, futures, and related options.
\end{principle}

\begin{caveat} \hlt{Post-Trade Complexity}\\
Options fragmentation is constrained by clearing, exercise, assignment, margin, and position-risk management. The trading interface may look electronic, but the economics are dominated by cross-contract risk.
\end{caveat}

\begin{mechanism} \hlt{Market-Maker-Centric Liquidity}\\
Options markets rely heavily on professional market makers posting two-sided liquidity. Buy-side traders typically consume liquidity because maintaining coherent quotes across many strikes and maturities requires specialised models and hedging infrastructure.
\end{mechanism}

\begin{taxonomy} \hlt{Options Liquidity Segmentation}
\begin{enumerate}[label=\roman*.]
\setlength{\itemsep}{0pt}
\item \hlt{moneyness}: at-the-money contracts are usually deepest;
\item \hlt{expiry}: near expiries trade differently from long-dated optionality;
\item \hlt{underlying liquidity}: hedging cost transfers into option spreads;
\item \hlt{volatility regime}: high uncertainty widens quotes and reduces size;
\item \hlt{complex orders}: spreads and multi-leg trades may be more relevant than outrights.
\end{enumerate}
\end{taxonomy}
